\newpage
\section{Generation Prompts and Details}
\label{sec:appendix:generationdetails}

Below we detail the structure of our prompts and inference parameters used for synthetic data generation.

\subsection{Extracting Topics from C4}
\label{sec:extracttopic}

We select random sentences from C4 \citep{t5andc4}, and extract the fine-grained topic of the sentence using GPT-4 \citep{gpt4} and the zero-shot prompt shown below. We perform sampling with \texttt{temperature} = 1.0 and \texttt{top\_p} = 0.0.

\vspace{1em}
{\small
\begin{verbatim}
 What is the fine-grained topic of the following text: {sentence} Only return the topic.    
\end{verbatim}}
\vspace{1em}

\noindent The fine-grained topic is then used as part of the attributed prompt described in Section \ref{sec:genposandneg} to ensure diversity in the generations.

\subsection{Generating Positive and Negative Example Sentences for Each Style Feature}
\label{sec:genposandneg}

For each style feature, we generate positive and negative parallel examples using a zero-shot prompt and the attributed prompt (AttrPrompt) method \citep{attrprompt} with GPT-4 to create diverse and realistic synthetic examples. To generate many examples per style feature, we use a new unused topic extracted from C4 in the prompt each time, we randomly sample a new permutation of the attributes in the prompt, and we perform sampling with \texttt{temperature} = 1.0 and \texttt{top\_p} = 1.0. We demonstrate an example below for the ``Usage of Active Voice'' style feature.

\vspace{1em}
{\small
\begin{verbatim}
Generate a pair of active and passive sentences with the following attributes:
    1. Topic: {topic}
    2. Length: {sentence_length}
    3. Point of view: {point_of_view}
    4. Tense: {tense}
    5. Type of Sentence: {sentence_type}
    
Ensure that the generated sentences meet the following conditions:
    1. There is no extra information in one sentence that is not in the other. 
    2. The difference between the two sentences is subtle. 
    3. The two sentences have the same length.
    {special_conditions_for_style_feature}
Use Format:
    Active: [sentence]
    Passive: [sentence]
Your response should only consist of the two sentences, without quotation marks.
\end{verbatim}}
\vspace{1em}

\noindent For the exact prompts for each style feature, see the code in our supplementary materials for this work.